% !TEX root = ../../main.tex
\section{Các chỉ tiêu đánh giá chất lượng điều khiển}
\label{sec:background-metrics}

\subsection{Sai lệch bám và dao động trong trạng thái xác lập}

Chất lượng điều khiển được đánh giá từ đáp ứng đo trên mô hình vật lý. Các chỉ tiêu được tính riêng cho roll và pitch, với cùng quy ước đơn vị và khoảng thời gian khảo sát. Ký hiệu $\chi$ đại diện cho một trong hai góc $\phi$ hoặc $\theta$, $\chi_d[k]$ là góc đặt và $\hat\chi[k]$ là góc ước lượng tại mẫu thứ $k$. Sai lệch bám được xác định bởi:
\begin{equation}
    e_\chi[k]=\chi_d[k]-\hat\chi[k].
    \label{eq:background-angle-error}
\end{equation}
Góc trong các phương trình điều khiển sử dụng radian; khi biểu diễn bằng độ trên đồ thị, cả góc đặt và góc ước lượng được chuyển đổi theo cùng hệ số. Việc ghi nhận thời điểm lấy mẫu cho phép liên hệ biến thiên góc với đáp ứng của lệnh động cơ.

Trên một tập gồm $N$ mẫu, căn quân phương của sai lệch (RMSE) và sai lệch tuyệt đối lớn nhất được tính theo:
\begin{equation}
    e_{\mathrm{RMS},\chi}=\sqrt{\frac{1}{N}\sum_{k=1}^{N}e_\chi^2[k]},
    \qquad e_{\max,\chi}=\max_{1\leq k\leq N}|e_\chi[k]|.
    \label{eq:background-error-indices}
\end{equation}

\clubpenalty=10000
RMSE phản ánh mức sai lệch trong toàn khoảng đánh giá, còn $e_{\max,\chi}$ cho biết độ lệch lớn nhất xuất hiện trong khoảng đó. Hai chỉ tiêu bổ sung cho nhau khi đáp ứng có những thời điểm lệch góc ngắn nhưng biên độ lớn.

Khi góc đặt không đổi, chọn một khoảng xác lập gồm $N_w$ mẫu để tính sai lệch trung bình và độ lệch chuẩn:
\begin{equation}
    \begin{aligned}
        \bar e_\chi&=\frac{1}{N_w}\sum_{k\in\mathcal W}e_\chi[k],\\
        \sigma_{e,\chi}&=\sqrt{\frac{1}{N_w}\sum_{k\in\mathcal W}
            \bigl(e_\chi[k]-\bar e_\chi\bigr)^2}.
    \end{aligned}
    \label{eq:background-steady-error}
\end{equation}
Ở đây, $\mathcal W$ là tập chỉ số thuộc khoảng xác lập. Giá trị $\bar e_\chi$ thể hiện sai lệch duy trì, còn $\sigma_{e,\chi}$ mô tả mức phân tán của tín hiệu phản hồi quanh giá trị trung bình. Dao động cơ học và nhiễu đo đều có thể ảnh hưởng đến đại lượng này, nên diễn biến góc được xem xét cùng vận tốc góc và lệnh động cơ.

\subsection{Độ quá điều chỉnh và thời gian xác lập}

Đáp ứng quá độ được khảo sát khi thay đổi góc đặt hoặc khi tư thế bị lệch khỏi trạng thái mong muốn. Độ quá điều chỉnh và thời gian xác lập là các chỉ tiêu thường dùng để mô tả mức vượt giá trị cuối và tốc độ ổn định của đáp ứng~\cite{astrom2021feedback}.

Với bước thay đổi góc đặt từ $\chi_{d0}$ sang $\chi_{d1}$ tại thời điểm $t_0$, đặt $\Delta\chi_d=\chi_{d1}-\chi_{d0}\ne0$ và $s_\chi=\operatorname{sgn}(\Delta\chi_d)$. Độ quá điều chỉnh theo phần trăm được xác định theo hướng của bước đặt:
\begin{equation}
    M_{p,\chi}=\frac{\max\!\left(0,\ \max_{k\in\mathcal K}
        \left[s_\chi\bigl(\hat\chi[k]-\chi_{d1}\bigr)\right]\right)}
        {|\Delta\chi_d|}\,100\%.
    \label{eq:background-overshoot}
\end{equation}
Trong đó, $\mathcal K$ là tập mẫu sau bước đặt và trước lần thay đổi lệnh tiếp theo. Cách xác định này áp dụng cho cả bước tăng và bước giảm góc đặt. Với phép thử khôi phục tư thế về góc đặt không đổi, độ lệch cực đại được biểu diễn trực tiếp bằng độ hoặc radian để tránh chuẩn hóa theo một bước đặt bằng không.

\begin{samepage}
Thời gian xác lập được tính từ thời điểm kích thích đến khi sai lệch đi vào dải cho phép và duy trì trong dải đó đến hết khoảng quan sát. Với chu kỳ lấy mẫu $T_s$, mẫu bắt đầu $k_0$ và mẫu kết thúc $k_f$, chọn $N_h\geq2$ là số mẫu duy trì tối thiểu để xác nhận đáp ứng. Tập chỉ số xét là $\mathcal K_s=\{k_0,\ldots,k_f-N_h+1\}$, từ đó:
\begin{equation}
    t_{s,\chi}=\min_{k\in\mathcal K_s}\left\{(k-k_0)T_s:\
        |e_\chi[j]|\leq\varepsilon_\chi,
        \ \forall j=k,\ldots,k_f\right\}.
    \label{eq:background-settling-time}
\end{equation}
\end{samepage}
Ngưỡng $\varepsilon_\chi$ được chọn phù hợp với yêu cầu giữ góc và độ phân giải của phép đo. Số mẫu $N_h$ xác định khoảng duy trì tối thiểu $(N_h-1)T_s$; cả hai thông số được sử dụng thống nhất giữa các lần thử. Khi không có chỉ số thỏa mãn hoặc dữ liệu sau lần vào dải chưa đủ $N_h$ mẫu, thời gian xác lập được ghi nhận là chưa xác định trong khoảng quan sát. Với chu kỳ lấy mẫu không đều, khoảng thời gian được tính từ dấu thời gian thực tế và điều kiện duy trì được kiểm tra theo thời gian.

\subsection{Khả năng khôi phục tư thế và điều kiện đánh giá}

Khả năng khôi phục tư thế được khảo sát bằng tác động gây lệch góc trong miền làm việc của hệ thống. Khi tác động kết thúc, góc đặt được giữ không đổi và theo dõi diễn biến góc trở về dải sai lệch cho phép. Các đại lượng dùng để đánh giá gồm độ lệch góc cực đại, thời gian khôi phục, mức dao động còn lại và sự xuất hiện bão hòa lệnh động cơ. Những đại lượng này phản ánh đồng thời khả năng hiệu chỉnh sai lệch và giới hạn tác động của hệ động lực.

Các phép thử được mô tả cùng cấu hình phần cứng, điều kiện nguồn điện, góc đặt, cách tạo sai lệch và khoảng thời gian ghi dữ liệu. Khi so sánh các bộ hệ số điều khiển, điều kiện thử và cách tính chỉ tiêu được giữ thống nhất.

Các chỉ tiêu bám tính theo $\hat\chi$ phản ánh chất lượng điều khiển đối với tín hiệu tư thế ước lượng; độ chính xác tuyệt đối của góc được xác định bằng phép đối chiếu với chuẩn đo góc độc lập. Phương án đo, điều kiện thực nghiệm và kết quả đánh giá của đề tài được trình bày trong \charef{chap:evaluation}.
